% ============================================================
%  Příprava na maturitu – Mechatronika
% ============================================================

\documentclass[a4paper,10pt]{article}

\usepackage[T1]{fontenc}
\usepackage[utf8]{inputenc}
\usepackage[czech]{babel}
\usepackage[left=2.4cm, right=1.8cm, top=1.8cm, bottom=1.8cm]{geometry}
\usepackage{lmodern}
\usepackage{microtype}
\usepackage{parskip}
\usepackage{titlesec}
\usepackage{xcolor}
\usepackage{hyperref}
\usepackage{tabularx}
\usepackage{booktabs}
\usepackage{tikz}
\usepackage{mdframed}
\usepackage{amsmath}

% --- Barvy ---
\definecolor{accent}{HTML}{03254E}
\definecolor{lightblue}{HTML}{EAF2F8}
\definecolor{lightyellow}{HTML}{FDFBE4}
\definecolor{lightgreen}{HTML}{EAF7EA}

% --- Boxy ---
\newmdenv[backgroundcolor=lightblue,linecolor=accent!40,roundcorner=4pt]{pojmybox}
\newmdenv[backgroundcolor=lightyellow,linecolor=accent!30,roundcorner=4pt]{vysvetlenibox}
\newmdenv[backgroundcolor=lightgreen,linecolor=accent!30,roundcorner=4pt]{vzorcebox}

\begin{document}

% ============================================================
\section{Senzorika -- poloha, rychlost, zrychlení, průtok, hladina}

\begin{pojmybox}
enkodér, LVDT, Hallův jev, MEMS, akcelerometr,
tachogenerátor, průtokoměr, hladinoměr,
citlivost, přesnost, rozlišení, linearita
\end{pojmybox}

% ============================================================
\begin{vysvetlenibox}

\textbf{Parametry senzorů}

\begin{vzorcebox}
\[
S = \frac{\Delta y}{\Delta x}
\]
\end{vzorcebox}

\begin{itemize}
\item \textbf{Citlivost} – změna výstupu na změnu vstupu
\item \textbf{Rozlišení} – nejmenší měřitelná změna
\item \textbf{Přesnost} – odchylka od skutečné hodnoty
\item \textbf{Linearita} – odchylka od ideální přímky
\end{itemize}

% ============================================================
\textbf{Snímače polohy}

\begin{tabularx}{\linewidth}{l X}
\toprule
Typ & Princip \\
\midrule
Inkrementální enkodér & počítá impulzy (relativní poloha) \\
Absolutní enkodér & každá poloha má unikátní kód \\
Potenciometr & změna odporu dle polohy \\
LVDT & změna indukčnosti (lineární posuv) \\
Hallův snímač & detekce magnetického pole \\
\bottomrule
\end{tabularx}

% ============================================================
\textbf{Rychlost a zrychlení}

\begin{vzorcebox}
\[
v = \frac{ds}{dt}, \quad a = \frac{dv}{dt}
\]
\end{vzorcebox}

\begin{itemize}
\item \textbf{Tachogenerátor} – napětí úměrné otáčkám:
\[
U = k \cdot n
\]

\item \textbf{Enkodér} – rychlost z impulzů:
\[
\omega = \frac{\Delta \varphi}{\Delta t}
\]

\item \textbf{Akcelerometr (MEMS)}
princip: malá hmotnost se při zrychlení posune → změna kapacity nebo odporu
\end{itemize}

% ============================================================
\textbf{Průtokoměry}

\begin{tabularx}{\linewidth}{l X}
\toprule
Typ & Princip \\
\midrule
Elektromagnetický & indukce napětí ve vodivé kapalině \\
Ultrazvukový & měření doby průchodu signálu \\
Vortex & vznik vírů za překážkou \\
Coriolis & měření hmotnostního průtoku (setrvačné síly) \\
\bottomrule
\end{tabularx}

% ============================================================
\textbf{Snímače hladiny}

\begin{tabularx}{\linewidth}{l X}
\toprule
Typ & Princip \\
\midrule
Plovákový & mechanický plovák kopíruje hladinu \\
Kapacitní & změna kapacity dle hladiny \\
Ultrazvukový & měření vzdálenosti od hladiny \\
Tlakový & hydrostatický tlak: $p = \rho g h$ \\
\bottomrule
\end{tabularx}

% ============================================================
\end{vysvetlenibox}

\end{document}
